\documentclass{article}
\usepackage[UTF8]{ctex}
\usepackage{amsmath,amssymb,amsfonts}
\usepackage{enumitem}
\usepackage{geometry}
\geometry{a4paper, margin=1in}
% === 错题本排版优化（不参与题目标识） ===
\setlength{\parindent}{0pt}
\setlength{\parskip}{0.35em plus 0.10em minus 0.08em}
\setlength{\emergencystretch}{3em}
\setlist{itemsep=0.25em, topsep=0.40em, parsep=0pt}
\allowdisplaybreaks[2]
% === 错题本排版优化结束 ===
\begin{document}
\section*{北京工业大学2024年高等代数真题}
\section*{6．北京工业大学2024年高等代数试卷}
一．设数域 \(P\) 上 \(n\) 阶方阵

\[
A=\left(\begin{array}{cccc}
a_{1}^{2} & a_{1} a_{2}+1 & \cdots & a_{1} a_{n}+1 \\
a_{2} a_{1}+1 & a_{2}^{2} & \cdots & a_{2} a_{n}+1 \\
\vdots & \vdots & & \vdots \\
a_{n} a_{1}+1 & a_{n} a_{2}+1 & \cdots & a_{n}^{2}
\end{array}\right) .
\]

其中 \(\sum_{i=1}^{n} a_{i}=1, \sum_{i=1}^{n} a_{i}^{2}=n\) 。

\begin{itemize}

  \item[1．] 求 \(A\) 的行列式 \(|A|\) 。

  \item[2．] 求 \(A\) 的所有特征值．
\end{itemize}

二．已知 \(V\) 是有理数域上的 3 维线性空间，\(\sigma\) 是 \(V\) 的线性变换，已知 \(\alpha, \beta, \gamma \in V\) 满足

\[
\sigma(\alpha)=\beta, \sigma(\beta)=\gamma, \sigma(\gamma)=\alpha-2 \beta .
\]

证明：若 \(\alpha \neq 0\) ，则 \(\alpha, \beta, \gamma\) 是 \(V\) 的一组基．

三．已知

\[
A=\left(\begin{array}{ccc}
1 & -1 & 1 \\
x & 4 & y \\
-1 & -1 & 3
\end{array}\right)
\]

有3个线性无关的特征向量，且 \(|A|=16\) ，其中 4 是 \(A\) 的一个特征值，求：

\begin{itemize}

  \item[1．] \(x, y\) 的值。

  \item[2．] \(A^{m}\) ，其中 \(m\) 是正整数。
\end{itemize}

四．解答如下问题：

\begin{itemize}

  \item[1．] 已知 \(A\) 为 \(n\) 阶半正定矩阵，\(B\) 是实矩阵且 \(A B+B A=O\) ，证明：\(A B=B A=O\) ．
\end{itemize}

2．已知 \(A, C\) 均为 \(n\) 阶正定矩阵，实矩阵 \(B\) 是 \(A X+X A=C\) 的唯一解，证明：\(B\) 为正定矩阵。

五．已知 \(V\) 是实数域上的 \(n\) 维线性空间，\(\sigma\) 是 \(V\) 的可对角化线性变换，\(U\) 是 \(V\) 的任一 \(\sigma\)－子空间．\\

1．证明：存在 \(V\) 的 \(\sigma\)－子空间 \(U^{\prime}\) ，使得 \(V=U \oplus U^{\prime}\) 。\\

2．证明：\(\sigma\) 在 \(U\) 上的限制线性变换 \(\left.\sigma\right|_{U}\) 也是可对角化的。

六．设 \(V\) 是实数域 \(\mathbb{R}\) 上所有 2 阶方阵构成的线性空间，\(\forall A, B \in V\) 定义 \((A, B)=\operatorname{tr}\left(A^{\prime} B\right)\) ．

\begin{itemize}

  \item[1．] 证明 \(V\) 是 4 维欧氏空间。

  \item[2．] 设 \(O \neq A \in V\) ，定义 \(\sigma_{A}(B)=B-\frac{(2 B, A)}{(A, A)} A\) ，证明 \(\sigma_{A}\) 是可对角化的线性变换．

  \item[3．] 设 \(A=\left(\begin{array}{ll}1 & 1 \\ 0 & 0\end{array}\right)\) ，求 \(V\) 的一组基，使得 \(\sigma_{A}\) 在这组基下的矩阵为对角矩阵，并求出该对角矩阵。
\end{itemize}

七．数域 \(P\) 上的 \(n\) 阶矩阵 \(A\) 在数域 \(P\) 上的全部特征值为 \(\lambda_{i}(i=1,2, \cdots, n)\) ．

\begin{itemize}

  \item[1．] 证明矩阵 \(A\) 相似于上三角矩阵 \(\left(\begin{array}{cccc}\lambda_{1} & * & \cdots & * \\ & \lambda_{2} & \ddots & \vdots \\ & & \ddots & * \\ & & & \lambda_{n}\end{array}\right)\) ．
\end{itemize}

2．设 \(\lambda_1,\ldots,\lambda_n\) 为 \(A\) 的全部特征值（按代数重数计）。若 \(-\lambda_{i}(i=1,2, \cdots, n)\) 全都不是 \(A\) 的特征值，证明 \(\varphi: x \mapsto x A+A^{\prime} x\) 是 \(M_{n}(P)\) 上的同构线性变换，其中 \(x \in M_{n}(P), M_{n}(P)\) 是数域 \(P\) 上所有 \(n\) 阶方阵做矩阵加法和矩阵数乘而构成的线性空间。
\end{document}
